% ==============================================================================
% 06_criteria3_teaching.tex — Criterion 3: Teaching and Learning Approach
% Score: 2
% ==============================================================================

\criterion{3}{แนวทางการเรียนการสอน}{Teaching and Learning Approach}

% ------------------------------------------------------------------------------
\criterionsub{3.1 The educational philosophy is shown to be articulated and communicated to all stakeholders. It is also shown to be reflected in the teaching and learning activities.}

มหาวิทยาลัยราชภัฏอุตรดิตถ์กำหนดปรัชญาการจัดการศึกษาไว้อย่างเป็นทางการในประกาศมหาวิทยาลัยราชภัฏอุตรดิตถ์ เรื่อง ปรัชญาการจัดการศึกษา ออกโดยรองศาสตราจารย์~ดร.สุภาวิณี สัตยาภรณ์ อธิการบดี เมื่อวันที่~25 สิงหาคม พ.ศ.~2566 โดยผ่านมติคณะกรรมการสภาวิชาการ ครั้งที่~6/2566 อ้างอิงมาตรา~31(1) แห่ง พ.ร.บ.มหาวิทยาลัยราชภัฏอุตรดิตถ์ พ.ศ.~2547 (อ้างอิง~\hyperlink{ref:AUNQA-3-1}{AUNQA-3-1})

ปรัชญาดังกล่าวกำหนดแนวคิดหลักคือ \textbf{``สร้างผลลัพธ์การเรียนรู้ด้วยประสบการณ์เชิงบูรณาการ'' (Integrated Experiential Learning --- IEL)} ซึ่งประกอบด้วยองค์ประกอบสำคัญสี่ด้าน ได้แก่ (1)~ผลลัพธ์การเรียนรู้ (Learning Outcomes) ครอบคลุม Knowledge, Skills, Ethics และ Character; (2)~ประสบการณ์เชิงบูรณาการ (Integrated Field Experiences) เพื่อพัฒนาตนเอง พัฒนาท้องถิ่น สู่การเรียนรู้ตลอดชีวิต; (3)~ผลลัพธ์การเรียนรู้เฉพาะด้าน (SLOs); และ (4)~ผลลัพธ์การเรียนรู้ทั่วไป (GLOs) ที่มุ่งสร้างพลเมืองดิจิทัลและพลเมืองเข้มแข็ง

ปรัชญา IEL สื่อสารไปยังผู้มีส่วนได้ส่วนเสียผ่านหลายช่องทาง ได้แก่ การประชาสัมพันธ์ผ่านเว็บไซต์มหาวิทยาลัย กิจกรรมปฐมนิเทศนักศึกษา (12 ตุลาคม 2568 โดยรองอธิการบดี รศ.ดร.อิสระ อินจันทร์ และคณบดีคณะเทคโนโลยีอุตสาหกรรม ผศ.ดร.ชัชพล เกษวิริยะกิจ) ซึ่งนักศึกษาได้รับทราบแนวทางการจัดการเรียนการสอนของมหาวิทยาลัยตั้งแต่วันแรกที่เข้าศึกษา

ปรัชญา IEL สะท้อนในกิจกรรมการเรียนการสอนของหลักสูตรอย่างเป็นรูปธรรม กล่าวคือ มคอ.3 ทุกวิชาระบุ IEL ไว้ในทุก CLO ของทุกวิชา ตารางที่~\ref{tbl:iel_course_activity} ยกตัวอย่างกิจกรรมที่สอดคล้องกับปรัชญา IEL ในรายวิชาที่จัดสอนปีการศึกษา~2568 (ไม่ครบทุกรายวิชา แสดงเฉพาะที่มีหลักฐานการดำเนินงานจริง)

\begin{table}[H]
  \centering
  \caption{ตัวอย่างกิจกรรมการเรียนการสอนที่สอดคล้องกับปรัชญา IEL ปีการศึกษา~2568}
  \label{tbl:iel_course_activity}
  \begingroup
  \footnotesize
  \setlength{\tabcolsep}{4pt}
  \renewcommand{\arraystretch}{1.35}
  \begin{tabularx}{\linewidth}{|>{\raggedright\arraybackslash}p{3.2cm}|>{\raggedright\arraybackslash}X|}
    \hline
    \rowcolor{olive!60!black}
    \multicolumn{1}{|c|}{\textcolor{white}{\textbf{รายวิชา}}} & \multicolumn{1}{c|}{\textcolor{white}{\textbf{กิจกรรมการเรียนการสอนที่สอดคล้องกับปรัชญาการศึกษา (IEL)}}} \\
    \hline
    7015101 เทคโนโลยีอุบัติใหม่ฯ &
    สัปดาห์ที่~10 จัด \textbf{IEL Community Workshop} นำโดย ผศ.ดร.พิศิษฐ์ นาคใจ (วิทยากรหลัก) โดยมีนักศึกษาร่วมสนับสนุนการอบรม: (ก)~อบรม ``การใช้งาน AI เพื่อการศึกษา'' ให้ครูโรงเรียนเทศบาล~3 แห่ง (เทศบาลท่าอิฐ, วัดหนองผา, วัดเกษมจิตตาราม) เมื่อ~26 มี.ค.~2569; (ข)~ถ่ายทอด AI \& IoT Smart Farming ให้สหกรณ์ออมทรัพย์ครูอุตรดิตถ์ 2 ครั้ง (4 และ 18 ต.ค.~2568) --- สะท้อน IEL ``ประสบการณ์เชิงบูรณาการเพื่อพัฒนาท้องถิ่น'' \\
    \hline
    4095101 AI \& ML &
    Workshop โครงงานย่อย พัฒนาโมเดล Deep Learning ทดสอบและปรับปรุง อภิปรายจริยธรรม AI และโครงงานปลายภาค (กิจกรรมตามที่ระบุใน มคอ.5 ภาค~2/2568) --- สะท้อน IEL ``Learning Outcomes ที่บูรณาการ Knowledge + Skills'' \\
    \hline
    4095102 Advanced Programming for ML &
    ฝึก Git/GitHub และ deploy โมเดล ML ตามที่ระบุใน มคอ.5 ภาค~2/2568 --- สะท้อน IEL ``ผลลัพธ์เฉพาะด้าน (SLO) ของวิชาชีพ ML Engineer'' \\
    \hline
    7015906 ระเบียบวิธีวิจัยฯ &
    การเขียนโครงร่างวิจัยเชื่อมโยงโจทย์จริงในที่ทำงาน/ท้องถิ่นของนักศึกษา --- สะท้อน IEL ``Learning to Learn เพื่อการเรียนรู้ตลอดชีวิต'' \\
    \hline
    กิจกรรมเสริมหลักสูตร &
    GenAI Short Courses 4 คอร์ส (Healthcare AI, Smart Life AI, Teacher AI, เลือกเสรี --- เม.ย.--พ.ค.~2569) เปิดให้บุคคลภายนอกเข้าร่วม ส่งเสริม Lifelong Learning ตาม IEL GLO ``พลเมืองดิจิทัล''; Special Class ``เทคโนโลยีตู้อบ'' โดย อ.พิภพ มณีจำนง (มรอ.วิทยาเขตน่าน) 29 ม.ค.~2569 --- สะท้อน IEL ``Integrated Field Experience'' ผ่านความร่วมมือข้ามวิทยาเขต \\
    \hline
  \end{tabularx}
  \endgroup
\end{table}

% ------------------------------------------------------------------------------
\criterionsub{3.2 The teaching and learning activities are shown to allow students to participate responsibly in the learning process.}

หลักสูตรออกแบบกิจกรรมการเรียนการสอนที่เปิดโอกาสให้นักศึกษามีส่วนร่วมอย่างมีความรับผิดชอบในกระบวนการเรียนรู้ของตนเองในหลายมิติ มิติแรกคือการเลือกเส้นทางการเรียนรู้ โดยนักศึกษาเลือกระหว่างแผน~1 (วิชาการ) และแผน~2 (วิชาชีพ) พร้อมเลือกวิชาเลือกจากคลัง~22 วิชาเพื่อออกแบบเส้นทางการศึกษาตามความสนใจและเป้าหมายอาชีพ มิติที่สองคือการกำหนดหัวข้อวิจัยด้วยตนเอง โดยนักศึกษาทั้ง~4 คนในรุ่นแรกต่างเสนอหัวข้อวิทยานิพนธ์/สารนิพนธ์ที่สอดคล้องกับความสนใจและบริบทการทำงาน ได้แก่ การตรวจจับผู้ขับขี่ไม่สวมหมวกด้วย Deep Learning (นายเกริกเกียรติ) ระบบ IoT สำหรับตู้อบข้าวแต๋น (นางสาวพอใจ) ระบบ AI สำหรับ EV charging station ของ PEA (นายจาตุรงค์) และการวิเคราะห์ภาพดาวเทียมสำหรับพื้นที่ปลูกอ้อย (นายจิรกิตติ์)

มิติที่สามคือกิจกรรม Active Participation ในชั้นเรียน โดยมคอ.3 วิชาบังคับทุกวิชาระบุกิจกรรมที่ให้นักศึกษาเป็นผู้ลงมือปฏิบัติและร่วมตัดสินใจในกระบวนการเรียนรู้ ตารางที่~\ref{tbl:student_participation_c32} ยกตัวอย่างกิจกรรมที่นักศึกษามีส่วนร่วมในการตัดสินใจ (ไม่ครบทุกรายวิชา แสดงเฉพาะที่มีหลักฐานการดำเนินงานจริงในปีการศึกษา~2568)

\begin{table}[H]
  \centering
  \caption{ตัวอย่างกิจกรรมการเรียนการสอนที่นักศึกษาได้มีส่วนร่วมในการตัดสินใจ ปีการศึกษา~2568}
  \label{tbl:student_participation_c32}
  \begingroup
  \footnotesize
  \setlength{\tabcolsep}{4pt}
  \renewcommand{\arraystretch}{1.35}
  \begin{tabularx}{\linewidth}{|>{\raggedright\arraybackslash}p{3.2cm}|>{\raggedright\arraybackslash}X|}
    \hline
    \rowcolor{olive!60!black}
    \multicolumn{1}{|c|}{\textcolor{white}{\textbf{รายวิชา}}} & \multicolumn{1}{c|}{\textcolor{white}{\textbf{กิจกรรมการเรียนการสอนที่นักศึกษาได้มีส่วนร่วมในการตัดสินใจ}}} \\
    \hline
    7015906 ระเบียบวิธีวิจัยฯ &
    นักศึกษาทั้ง~4 คนเลือกหัวข้อวิทยานิพนธ์/สารนิพนธ์ของตนเองจากบริบทการทำงาน: การตรวจจับผู้ขับขี่ไม่สวมหมวก (เกริกเกียรติ), IoT ตู้อบข้าวแต๋น (พอใจ), AI สำหรับ EV charging ของ PEA (จาตุรงค์), AI วิเคราะห์ภาพดาวเทียมพื้นที่อ้อย (จิรกิตติ์) --- รับผิดชอบขอบเขต ระเบียบวิธี และผลลัพธ์การวิจัยด้วยตนเอง \\
    \hline
    4095101 AI \& ML &
    โครงงานย่อยและโครงงานปลายภาค: นักศึกษาเลือกขอบเขตการพัฒนาโมเดล DL และนำเสนอผลงาน (ตามที่ระบุใน มคอ.5) \\
    \hline
    4095102 Advanced Programming for ML &
    โครงงาน deploy โมเดล ML: นักศึกษากำหนดขอบเขตและออกแบบ workflow ของโครงงานในรูปแบบกลุ่ม (ตามที่ระบุใน มคอ.5) \\
    \hline
    7015101 เทคโนโลยีอุบัติใหม่ฯ &
    IEL Community Workshop: นักศึกษาร่วมสนับสนุนการอบรม AI ให้กลุ่มเป้าหมายภายนอก ภายใต้การนำของ ผศ.ดร.พิศิษฐ์ นาคใจ (วิทยากรหลัก) ได้แก่ ครูเทศบาล 3 โรงเรียน (26 มี.ค.~2569) และสหกรณ์ออมทรัพย์ครูอุตรดิตถ์ (4 และ 18 ต.ค.~2568) โดยนักศึกษาช่วยเตรียมตัวอย่างและสาธิตเนื้อหาเฉพาะส่วน \\
    \hline
    การเลือกแผนการศึกษา &
    นักศึกษาทุกคนเลือกระหว่างแผน~1 (วิชาการ--วิทยานิพนธ์) หรือแผน~2 (วิชาชีพ--สารนิพนธ์) และเลือกวิชาเลือกจากคลัง~22 วิชาเพื่อกำหนดเส้นทางการเรียนรู้ของตน \\
    \hline
  \end{tabularx}
  \endgroup
\end{table}

นอกจากนี้ นักศึกษาระดับบัณฑิตศึกษายังได้ร่วมสนับสนุนการอบรม AI ให้ครูโรงเรียนเทศบาล~3 แห่ง (26 มีนาคม~2569) นำโดย ผศ.ดร.พิศิษฐ์ นาคใจ ซึ่งเป็นกิจกรรมที่ส่งเสริม Active Participation ผ่านประสบการณ์จริงนอกห้องเรียน

% ------------------------------------------------------------------------------
\criterionsub{3.3 The teaching and learning activities are shown to involve active learning by the students.}

การเรียนรู้เชิงรุก (Active Learning) ถูกกำหนดเป็นแกนกลางของการจัดการเรียนการสอน \textbf{ครบทุกรายวิชาทุกชั้นปี} โดยมีหลักฐานจาก Teaching Activities ที่ระบุในมคอ.3 ทุกวิชา ดังแสดงในตาราง~\ref{tbl:active_learning} (วิชาบังคับ~6 วิชา; วิชาเลือก~22 วิชาดู~\hyperlink{ref:AUNQA-3-2}{AUNQA-3-2})

\begin{table}[H]
  \centering
  \caption{กิจกรรมการเรียนการสอนแบบ Active Learning --- ทุกรายวิชาบังคับ ทุกชั้นปี}
  \label{tbl:active_learning}
  \begingroup
  \footnotesize
  \setlength{\tabcolsep}{3pt}
  \renewcommand{\arraystretch}{1.2}
  \begin{tabularx}{\linewidth}{|>{\raggedright\arraybackslash}p{3.2cm}|>{\raggedright\arraybackslash}X|>{\centering\arraybackslash}p{1.6cm}|}
    \hline
    \rowcolor{olive!60!black}
    \multicolumn{1}{|c|}{\textcolor{white}{\textbf{รายวิชา}}} & \multicolumn{1}{c|}{\textcolor{white}{\textbf{กิจกรรมการเรียนการสอนแบบ Active Learning}}} & \textcolor{white}{\textbf{CLO ที่ฝึก}} \\
    \hline
    4095101 AI and ML & พัฒนาโมเดล DL จริง, ทดสอบ, นำเสนอผลงาน, อภิปรายจริยธรรม AI & CLO3,4,5 \\
    \hline
    4095102 Advanced Programming for ML & ฝึก Git/GitHub, deploy โมเดลใน production, โครงงานวิเคราะห์ข้อมูลจริง & CLO2,3,5 \\
    \hline
    7015101 Emerging Tech & IEL community workshop สัปดาห์~10, ทดลอง IoT/Blockchain จริง & CLO2,4,5 \\
    \hline
    7015906 Research Methodology & วิเคราะห์ปัญหาวิจัย, ออกแบบ methodology, เขียนโครงร่างวิจัย & CLO2,3,5 \\
    \hline
    7015907 Seminar (เปิดสอนชั้นปี~2 --- กิจกรรมตามแผน มคอ.3) & นำเสนอบทความวิจัย, วิพากษ์งานวิจัยเพื่อน, อภิปรายกลุ่ม & CLO1,2,3 \\
    \hline
    7015908 หัวข้อพิเศษฯ (เปิดสอนชั้นปี~2 --- กิจกรรมตามแผน มคอ.3) & Problem-based Learning, วิเคราะห์เทคโนโลยี, พัฒนานวัตกรรม, Portfolio & CLO2,3,5,6 \\
    \hline
  \end{tabularx}
  \endgroup
\end{table}

วิชาเลือก~22 วิชา (กลุ่ม AI, IoT, Networks, Software Engineering) ทุกวิชามีตาราง Teaching \& Learning Activities ระบุใน มคอ.3 ครบถ้วน รายละเอียดปรากฏใน evidence index (อ้างอิง~\hyperlink{ref:AUNQA-3-2}{AUNQA-3-2})

นอกจากนี้ หลักสูตรได้จัด Special Class โดยเชิญวิทยากรพิเศษจากภายนอก ซึ่งเปิดโอกาสให้นักศึกษาเรียนรู้จากประสบการณ์จริงของผู้ประกอบวิชาชีพ ได้แก่ อ.ธีรศักดิ์ อุปการ บรรยายพิเศษหัวข้อ ``โดรน AI สำหรับตรวจสอบหลังคาอาคาร'' (12 มกราคม 2569) อ.คชรัตน์ ภูฆัง และ ผศ.ดร.กชกร ลาภมาก บรรยาย Backward Design สำหรับการออกแบบหลักสูตร (19 มกราคม 2569) และ อ.พิภพ มณีจำนง (มรอ. วิทยาเขตน่าน) บรรยาย Special Class เทคโนโลยีตู้อบ (29 มกราคม 2569) ซึ่งกิจกรรมเหล่านี้ส่งเสริมให้นักศึกษาเชื่อมโยงทฤษฎีกับการปฏิบัติจริง

% ------------------------------------------------------------------------------
\criterionsub{3.4 The teaching and learning activities are shown to promote learning, learning how to learn, and instilling in students a commitment for life-long learning (e.g., commitment to critical inquiry, information-processing skills, and a willingness to experiment with new ideas and practices).}

หลักสูตร CPE\&AI กำหนดกรอบ LLL ใน 3 มิติที่วัดได้ ได้แก่ (1) ทักษะการสืบค้นและประมวลผลความรู้วิชาการด้วยตนเอง (critical inquiry) (2) ความสามารถเรียนรู้เครื่องมือและเทคโนโลยีใหม่อย่างอิสระ (self-directed technology learning) และ (3) การทดลองประยุกต์แนวคิดใหม่กับโจทย์จริง (experimentation with new ideas) ซึ่งสอดคล้องกับลักษณะของสาขาวิชาที่ภูมิทัศน์เทคโนโลยีเปลี่ยนแปลงรวดเร็ว

การส่งเสริม LLL ในหลักสูตรนี้ดำเนินการผ่านการบูรณาการใน CLO รายวิชา กิจกรรมการสอน และโครงการขยายความรู้สู่ชุมชน ซึ่งปรากฏชัดเจนในสามระดับ

ระดับรายวิชา มคอ.3 วิชา 7015906 ระเบียบวิธีวิจัยฯ กำหนด CLO5 ``ประยุกต์ใช้วิธีการทางสถิติในการวิเคราะห์ข้อมูลและสรุปผลการวิจัย'' ซึ่ง map กับ PLO5 (คิดวิเคราะห์ การสื่อสาร การทำงานร่วมกัน การปรับตัว และการเรียนรู้ตลอดชีวิต) ทั้งนี้ วิชา 7015101 เทคโนโลยีอุบัติใหม่ฯ ยังกำหนด CLO4 ``ตอบสนองต่อการเรียนรู้เทคโนโลยีใหม่ด้วยความกระตือรือร้นและมีส่วนร่วม'' ซึ่งเป็น Affective objective ที่ปลูกฝัง disposition ต่อการเรียนรู้ตลอดชีวิต

ระดับหลักสูตร วิชาเสริม 2 วิชา ได้แก่ 1555101 ภาษาอังกฤษสำหรับนักศึกษาบัณฑิตศึกษา และ 1506001 การเขียนวิทยานิพนธ์หรือสารนิพนธ์ในระดับบัณฑิตศึกษา เสริมสร้างทักษะพื้นฐานสำหรับการเรียนรู้ตลอดชีวิต ทั้งทักษะภาษาอังกฤษสำหรับเข้าถึงองค์ความรู้สากล และทักษะการสื่อสารวิชาการ

ระดับมหาวิทยาลัยและชุมชน หลักสูตรจัดโปรแกรม GenAI Short Courses ที่เปิดให้บุคคลภายนอก~4 คอร์ส (เมษายน--พฤษภาคม 2569) ซึ่งสะท้อนวัฒนธรรม LLL ที่หลักสูตรสนับสนุนทั้งในและนอกองค์กร (อ้างอิง~\hyperlink{ref:AUNQA-C6-STSAT-2568}{AUNQA-C6-STSAT-2568})

% ------------------------------------------------------------------------------
\criterionsub{3.5 The teaching and learning activities are shown to inculcate in students, new ideas, creative thought, innovation, and an entrepreneurial mindset.}

ความคิดสร้างสรรค์ นวัตกรรม และ Entrepreneurial Mindset ถือเป็นจุดมุ่งหมายหนึ่งของการจัดการเรียนการสอนในหลักสูตร โดยปรากฏอย่างเป็นรูปธรรมในกิจกรรมหลายรูปแบบ ดังนี้

ในด้านความคิดริเริ่มและความคิดสร้างสรรค์ วิชา 7015101 ที่นำโจทย์จริงจากชุมชนมาเป็นโจทย์การเรียนรู้ผ่าน IEL Community Workshop กระตุ้นให้นักศึกษาคิดหาแนวทางแก้ปัญหาเชิงสร้างสรรค์ที่ไม่มีคำตอบตายตัว สำหรับวิชา 7015907 สัมมนา และ 7015908 หัวข้อพิเศษฯ มีกำหนดเปิดสอนในชั้นปีที่~2 ตามแผนการศึกษา

สำหรับนวัตกรรม วิทยานิพนธ์และสารนิพนธ์ของนักศึกษาทุกคนมุ่งพัฒนาต้นแบบที่แก้ปัญหาจริง ได้แก่ การตรวจจับผู้ขับขี่ไม่สวมหมวกด้วย Deep Learning, ระบบ IoT สำหรับตู้อบข้าวแต๋น, ระบบ AI สำหรับ EV charging station ของ PEA และการวิเคราะห์ภาพดาวเทียมสำหรับพื้นที่ปลูกอ้อย ผลงานเหล่านี้สะท้อนว่าหลักสูตรมุ่งเน้นนวัตกรรมที่มีประโยชน์ใช้งานได้จริงมากกว่าการเรียนรู้เชิงทฤษฎีเพียงอย่างเดียว

ในด้าน Entrepreneurial Mindset หลักสูตรจัดโปรแกรม GenAI Short Courses~4 คอร์ส (Healthcare AI, Smart Life AI, Teacher AI และคอร์สเลือกเสรี) ในเดือนเมษายน--พฤษภาคม~2569 โดยเปิดให้ผู้สนใจทั่วไปเข้าร่วม และนักศึกษาระดับบัณฑิตศึกษาได้ร่วมออกแบบและดำเนินการ กิจกรรมนี้ปลูกฝัง mindset การสร้างคุณค่าและการแบ่งปันความรู้สู่สังคมอย่างเป็นรูปธรรม

ตารางที่~\ref{tbl:creative_innovation_c35} สรุปกิจกรรมที่ส่งเสริม~4 ด้านตามเกณฑ์ 3.5 ครบทุกด้าน (แสดงเฉพาะรายวิชา/กิจกรรมที่มีหลักฐานการดำเนินงานจริงในปีการศึกษา~2568)

\begin{table}[H]
  \centering
  \caption{กิจกรรมที่ส่งเสริม New Ideas, Creative Thought, Innovation และ Entrepreneurial Mindset}
  \label{tbl:creative_innovation_c35}
  \begingroup
  \footnotesize
  \setlength{\tabcolsep}{3pt}
  \renewcommand{\arraystretch}{1.3}
  \begin{tabularx}{\linewidth}{|>{\raggedright\arraybackslash}p{2.4cm}|>{\raggedright\arraybackslash}X|>{\centering\arraybackslash}p{1.0cm}|>{\centering\arraybackslash}p{1.2cm}|>{\centering\arraybackslash}p{1.3cm}|>{\centering\arraybackslash}p{1.5cm}|}
    \hline
    \rowcolor{olive!60!black}
    \makecell[c]{\color{white}\bfseries รายวิชา/\\ \color{white}\bfseries โครงการ} &
    \makecell[c]{\color{white}\bfseries กิจกรรม} &
    \makecell[c]{\color{white}\bfseries New\\ \color{white}\bfseries Ideas} &
    \makecell[c]{\color{white}\bfseries Creative\\ \color{white}\bfseries Thought} &
    \makecell[c]{\color{white}\bfseries Innovation} &
    \makecell[c]{\color{white}\bfseries Entrepre-\\ \color{white}\bfseries neurial} \\
    \hline
    7015101 เทคโนโลยีอุบัติใหม่ฯ &
    IEL Community Workshop --- ผศ.ดร.พิศิษฐ์ นาคใจ เป็นวิทยากรหลัก นักศึกษาร่วมสนับสนุนการอบรม AI ให้ครูเทศบาล 3 โรงเรียน (26 มี.ค.~2569) และสหกรณ์ครูอุตรดิตถ์ (4, 18 ต.ค.~2568) &
    \checkmark & \checkmark & & \checkmark \\
    \hline
    วิทยานิพนธ์/สารนิพนธ์ (7015901--05) &
    นักศึกษาทั้ง~4 คนพัฒนาต้นแบบจากโจทย์จริง: ตรวจจับผู้ขับขี่ไม่สวมหมวก (DL), IoT ตู้อบข้าวแต๋น, AI สำหรับ EV charging ของ PEA, AI วิเคราะห์ภาพดาวเทียมพื้นที่อ้อย &
    \checkmark & \checkmark & \checkmark & \\
    \hline
    GenAI Short Courses &
    4 คอร์ส (Healthcare AI / Smart Life AI / Teacher AI / เลือกเสรี --- เม.ย.--พ.ค.~2569) นักศึกษาบัณฑิตศึกษาร่วมออกแบบและดำเนินการ เปิดให้บุคคลภายนอกเข้าร่วม &
    \checkmark & & \checkmark & \checkmark \\
    \hline
  \end{tabularx}
  \endgroup
  \begin{flushleft}
    \footnotesize หมายเหตุ: ไม่จำเป็นต้องครบทุกรายวิชา แต่กิจกรรมโดยรวมครอบคลุมครบทั้ง~4 ด้าน
  \end{flushleft}
\end{table}

% ------------------------------------------------------------------------------
\criterionsub{3.6 The teaching and learning processes are shown to be continuously improved to ensure their relevance to the needs of industry and are aligned to the expected learning outcomes.}

หลักสูตรวางระบบทบทวนและปรับปรุงกิจกรรมการเรียนการสอนอย่างต่อเนื่องผ่านกลไกสองระดับที่เชื่อมกัน ในระดับรายวิชา อาจารย์ผู้สอนทุกวิชาจัดทำ มคอ.5 เมื่อสิ้นสุดแต่ละภาคการศึกษา โดยรายงานผลการประเมิน วิเคราะห์ปัญหาที่พบ และเสนอแนวทางปรับปรุงกิจกรรมการเรียนการสอนสำหรับภาคการศึกษาถัดไป ในระดับหลักสูตร คณะกรรมการบริหารหลักสูตรประชุมทบทวนผล มคอ.5 ทุกภาคการศึกษา และพิจารณาปรับปรุงแผนการเรียนการสอนให้สอดคล้องกับความต้องการของภาคอุตสาหกรรมและแนวโน้มเทคโนโลยี การทบทวนโดยบุคลากรภายนอกหลักสูตร (Third-party Input) ดำเนินการผ่านหลายช่องทางในปีการศึกษา~2568 ได้แก่ การบรรยาย Backward Design โดย อ.คชรัตน์ ภูฆัง (คณะครุศาสตร์) และ ผศ.ดร.กชกร ลาภมาก (คณะวิทยาศาสตร์) เมื่อ~19 มกราคม~2569 เพื่อรับมุมมองจากผู้เชี่ยวชาญด้านการออกแบบการสอนจากสาขาอื่น Feedback จาก อ.ธีรศักดิ์ อุปการ และ อ.พิภพ มณีจำนง (มรอ.วิทยาเขตน่าน) ที่ยืนยันว่าเนื้อหา TLA ตรงกับความต้องการของภาคอุตสาหกรรมและบริบทข้ามวิทยาเขต ส่วน Formal Third-party Review อย่างเป็นทางการดำเนินการรอบแรกในการประชุมทบทวนผลสัมฤทธิ์ผู้เรียน ภาคการศึกษาที่~2/2568 เมื่อวันที่~27 มีนาคม~2569 (\hyperlink{ref:AUNQA-4-5-1}{AUNQA-4-5-1})

มคอ.5 ของ~4 วิชาบังคับ (4095101 และ 4095102 ในภาค~2/2568; 7015101 และ 7015906 ในภาค~3/2568) ทั้ง~4 ฉบับลงนามรายงานเมื่อ~30 เมษายน~2569 ระบุปัญหาและแผนการปรับปรุงที่จะนำไปใช้จริงในภาคถัดไป เป็นหลักฐาน PDCA Loop ที่สมบูรณ์ วิชา 4095101 AI and ML รายงานข้อจำกัดด้านการเข้าถึง GPU สำหรับเทรนโมเดลขนาดใหญ่และความหลากหลายของพื้นฐานคณิตศาสตร์ของนักศึกษา จึงมีแผนเพิ่มชุดแบบฝึกหัด pre-test เพื่อจัดกลุ่มช่วยเหลือ และจัดทำตัวอย่างโค้ดและ Notebook มาตรฐาน วิชา 4095102 Advanced Programming for ML รายงานปัญหาเรื่องการ deploy โมเดลใน production environment และความแตกต่าง of พื้นฐาน Python ของนักศึกษา จึงมีแผนจัดทำ Repository ตัวอย่างพร้อม CI สำหรับโครงงาน และเพิ่ม Workshop การใช้ MLOps เบื้องต้น วิชา 7015906 Research Methodology รายงานว่านักศึกษาบางกลุ่มต้องการเวลาฝึกซอฟต์แวร์สถิติเพิ่ม จึงมีแผนจัดทำตัวอย่างโครงร่างทบทวนวรรณกรรมฉบับมาตรฐาน และเพิ่ม Workshop การใช้เครื่องมือสถิติ (Python/R) ส่วนวิชา 7015101 Emerging Technology รายงานข้อจำกัดเรื่องจำนวนนักศึกษาน้อย ($n=4$) และการประสานงานสำหรับ IEL Community Workshop จึงมีแผนเพิ่มชุดข้อมูลฝึกปฏิบัติและตัวอย่างโครงงานจากภาคอุตสาหกรรม จัดทำแบบฟอร์มสะท้อนคิดมาตรฐานสำหรับ CLO3--CLO4 และเชื่อมโยงผล มคอ.5 กับการทบทวนระดับหลักสูตรในภาคเรียนถัดไป ผลลัพธ์จากกิจกรรม Community Engagement เป็นหลักฐานเสริมที่สะท้อนคุณภาพ TLA โดย ผศ.ดร.พิศิษฐ์ นาคใจ นำอบรม AI ให้ครูโรงเรียนเทศบาล~3 แห่ง (26 มีนาคม~2569) และนักศึกษาร่วมสนับสนุนกระบวนการอบรมจริง และนักศึกษารหัส~68 ของหลักสูตร (คุณจิรกิตติ์ วราภรณ์ ครูผู้ช่วยแผนกช่างอิเล็กทรอนิกส์ วิทยาลัยเทคนิคอุตรดิตถ์) ได้รับรางวัลชนะเลิศการประกวดสิ่งประดิษฐ์ เทคโนโลยีและนวัตกรรมเพื่อสุขภาพ งานวันเทคโนโลยีครั้งที่~11 (12 กุมภาพันธ์~2569) ยืนยันว่ากิจกรรมการเรียนการสอนสอดคล้องกับ CLOs และ PLOs อย่างมีประสิทธิผลในเชิงปฏิบัติ

\begin{strengthbox}
หลักสูตรมีปรัชญา IEL ที่กำหนดอย่างเป็นทางการโดยอธิการบดี (25 ส.ค. 2566) สะท้อนใน มคอ.3 ทุกวิชา มีกิจกรรม Active Learning ครอบคลุมทุกวิชาบังคับ นักศึกษาเลือกหัวข้อวิจัยเอง มี Community Engagement อบรม AI ให้โรงเรียน~3 แห่งและสหกรณ์ออมทรัพย์ครูอุตรดิตถ์ นำโดย ผศ.ดร.พิศิษฐ์ นาคใจ GenAI Short Courses~4 คอร์ส และ Cross-faculty input จากบุคลากรสาขาอื่น (ครุศาสตร์/วิทยาศาสตร์/วิทยาเขตน่าน) ผลสำรวจความพึงพอใจนักศึกษา $n=4$ ปีการศึกษา~2568 ให้คะแนนด้านคุณภาพการเรียนการสอน~5.00/5.00 (สูงสุด) สะท้อนประสิทธิผลของ TLA จริง
\end{strengthbox}

\begin{weaknessbox}
R3.6 Internal Third-party Review (ทวนสอบโดยคณะกรรมการบริหารหลักสูตรที่ไม่ได้สอนวิชานั้น) ดำเนินการรอบแรกแล้วเมื่อ~27 มี.ค.~2569 แต่ยังอยู่ในรูปรายงานประชุมภายใน ยังไม่เป็นแบบฟอร์มมาตรฐานรายวิชา; External Third-party Review โดยผู้ประเมินภายนอกสถาบันยังรอรอบการประกันคุณภาพ ปีการศึกษา~2568 สิ้นสุด; กิจกรรม LLL ต้องการกรอบนิยามที่ชัดเจนขึ้นในบริบทของหลักสูตร CPE\&AI เพื่อให้วัดได้และประเมินได้
\end{weaknessbox}

\begin{selfassessmentbox}{2}{หลักสูตรมีปรัชญา IEL ประกาศอย่างเป็นทางการโดยอธิการบดี สะท้อนใน มคอ.3 ทุกวิชา มีกิจกรรม Active Learning ครอบคลุมทุกวิชาบังคับ GenAI Short Courses~4 คอร์ส Cross-faculty input ทบทวน TLA และ PDCA Loop ผ่าน มคอ.5 ทุกวิชา ผลสำรวจนักศึกษา $n=4$ ให้คะแนนคุณภาพการสอน~5.00/5.00 Formal Third-party Review ดำเนินการรอบแรกแล้ว (27~มี.ค.~2569) แต่ยังไม่มีรูปแบบมาตรฐานเป็นแบบฟอร์มรายวิชา}
\end{selfassessmentbox}

\begin{evidencelist}
\evidence{AUNQA-3-1}{ประกาศ มรอ. เรื่องปรัชญาการจัดการศึกษา IEL (25 ส.ค. 2566) โดย รศ.ดร.สุภาวิณี สัตยาภรณ์}{\href{https://syweb.kidcbc.work/Eviden_aunqa68/\%E0\%B8\%A3\%E0\%B8\%B0\%E0\%B9\%80\%E0\%B8\%9A\%E0\%B8\%B5\%E0\%B8\%A2\%E0\%B8\%9A\%E0\%B8\%82\%E0\%B9\%89\%E0\%B8\%AD\%E0\%B8\%9A\%E0\%B8\%B1\%E0\%B8\%87\%E0\%B8\%84\%E0\%B8\%B1\%E0\%B8\%9A\%E0\%B9\%81\%E0\%B8\%A5\%E0\%B8\%B0\%E0\%B8\%9B\%E0\%B8\%A3\%E0\%B8\%B0\%E0\%B8\%81\%E0\%B8\%B2\%E0\%B8\%A8\%E0\%B8\%AD\%E0\%B8\%B7\%E0\%B9\%88\%E0\%B8\%99\%E0\%B9\%86/AUNQA-3-1_\%E0\%B8\%9B\%E0\%B8\%A3\%E0\%B8\%B0\%E0\%B8\%81\%E0\%B8\%B2\%E0\%B8\%A8IEL_2566.pdf}{\texttt{[PDF]}}}
\evidence{AUNQA-3-2}{มคอ.3 ทุกรายวิชา (ระบุ IEL ในทุก CLO)}{\href{https://syweb.kidcbc.work/Eviden_aunqa68/\%E0\%B8\%A1\%E0\%B8\%84\%E0\%B8\%AD3/AUN3_4095101.pdf}{\texttt{[PDF]}}}
\evidence{AUNQA-3-3}{สไลด์ปฐมนิเทศนักศึกษา CPE\&AI รุ่นแรก (12 ต.ค. 2568) โดย รองอธิการบดีและคณบดี}{\href{https://syweb.kidcbc.work/Eviden_aunqa68/\%E0\%B8\%A3\%E0\%B8\%B0\%E0\%B9\%80\%E0\%B8\%9A\%E0\%B8\%B5\%E0\%B8\%A2\%E0\%B8\%9A\%E0\%B8\%82\%E0\%B9\%89\%E0\%B8\%AD\%E0\%B8\%9A\%E0\%B8\%B1\%E0\%B8\%87\%E0\%B8\%84\%E0\%B8\%B1\%E0\%B8\%9A\%E0\%B9\%81\%E0\%B8\%A5\%E0\%B8\%B0\%E0\%B8\%9B\%E0\%B8\%A3\%E0\%B8\%B0\%E0\%B8\%81\%E0\%B8\%B2\%E0\%B8\%A8\%E0\%B8\%AD\%E0\%B8\%B7\%E0\%B9\%88\%E0\%B8\%99\%E0\%B9\%86/AUNQA-1-5_\%E0\%B8\%9B\%E0\%B8\%90\%E0\%B8\%A1\%E0\%B8\%99\%E0\%B8\%B4\%E0\%B9\%80\%E0\%B8\%97\%E0\%B8\%A8\%E0\%B8\%99\%E0\%B8\%B1\%E0\%B8\%81\%E0\%B8\%A8\%E0\%B8\%B6\%E0\%B8\%81\%E0\%B8\%A9\%E0\%B8\%B2\%E0\%B8\%9A\%E0\%B8\%B1\%E0\%B8\%93\%E0\%B8\%91\%E0\%B8\%B4\%E0\%B8\%95\%E0\%B8\%A8\%E0\%B8\%B6\%E0\%B8\%81\%E0\%B8\%A9\%E0\%B8\%B2_2568.pdf}{\texttt{[PDF]}}}
\evidence{AUNQA-C6-STSAT-2568}{แบบสำรวจความพึงพอใจนักศึกษาและผู้มีส่วนได้ส่วนเสีย CPE\&AI ปีการศึกษา~2568 ($n=4$): คะแนนด้านคุณภาพการสอน/ประเมิน~5.00/5.00 --- student\_satisfaction\_survey\_2568.csv}{\href{https://docs.google.com/forms/d/1xpfb6AfYkbOkidJNh6ci95ux4ZLYzjvpUuqARYjhr5k/viewanalytics}{\texttt{[Google Report]}}}
\end{evidencelist}
